\section{The harness}
\label{sec:platform-harness}

The platform lives in \texttt{Benchmarking/}. Its core is small: seven
modules of roughly 700 lines in \texttt{core/}, of which the LASSO metrics
module is a third; one adapter file per method family in
\texttt{adapters/}; and a command-line entry point. Everything else
(the sweep driver, the data generator, the experiment scripts) is a consumer
of this core. Figure~\ref{fig:platform-architecture} shows how the pieces fit.

% Figure 3.1 -- platform architecture.
% Requires in the preamble:
%   \usepackage{tikz}
%   \usetikzlibrary{positioning, arrows.meta, fit, backgrounds, calc}
\begin{figure}[htbp]
\centering
\begin{tikzpicture}[
  font=\small,
  box/.style={draw, rounded corners=2pt, minimum height=8mm, minimum width=26mm, align=center, fill=white},
  lib/.style={box, minimum width=22mm, minimum height=10mm, fill=gray!10},
  grp/.style={draw, dashed, rounded corners=4pt, inner sep=5pt},
  arr/.style={-{Latex[length=2mm]}, thick},
  lbl/.style={font=\footnotesize\itshape, gray!80!black},
  node distance=6mm and 8mm,
]
% --- consumers (top) ---
\node[box] (exp) {experiment script\\\texttt{experiments/*.py}};
\node[box, right=of exp] (sweep) {sweep driver\\\texttt{sweep.py}};
\node[box, right=of sweep] (cli) {CLI / suites\\\texttt{run.py}};

% --- core (middle) ---
\node[box, below=10mm of sweep] (runner) {runner\\\texttt{train\_evaluate(config)}};
\node[box, left=of runner] (registry) {registry\\\texttt{get\_method(name)}};
\node[box, right=of runner] (result) {\texttt{BenchmarkResult}\\\texttt{result.json}, trajectory};

% --- adapters ---
\node[box, below=10mm of runner, minimum width=70mm] (adapter)
  {adapter: \texttt{L2OMethod.train / evaluate}\\
   builds flags $\to$ \texttt{shell.run\_script} (validates flags, logs command) $\to$ reads artifacts};

% --- libraries (bottom) ---
\node[lib, below=10mm of adapter, xshift=-46mm] (dm) {L2O-DM /\\RNNProp};
\node[lib, right=3mm of dm] (scale) {L2O-Scale /\\Entropy / Jacobian};
\node[lib, right=3mm of scale] (swarm) {L2O-Swarm};
\node[lib, right=3mm of swarm] (mb) {LISTA /\\ALISTA};
\node[lib, right=3mm of mb] (cls) {Adam / SGD /\\Adagrad (host)};

% --- shared inputs (left) ---
\node[box, left=14mm of adapter, fill=gray!10, minimum width=24mm] (data)
  {seeded dataset\\\texttt{data/lasso.py}\\$A, b, x_\mathrm{true}, x_0$};
\node[box, below=4mm of data, fill=gray!10, minimum width=24mm] (metrics)
  {metrics\\\texttt{lasso\_metrics.py}\\$x^\ast$ by FISTA};

% --- groups ---
\begin{scope}[on background layer]
  \node[grp, fit=(exp)(sweep)(cli), label={[lbl]above left:consumers}] {};
  \node[grp, fit=(registry)(runner)(result)(adapter), label={[lbl]above left:\texttt{Benchmarking/core} + \texttt{adapters}}] {};
  \node[grp, fit=(dm)(scale)(swarm)(mb)(cls), label={[lbl]below left:migrated Open-L2O libraries, each run as a subprocess in its own directory}] {};
\end{scope}

% --- arrows ---
\draw[arr] (exp) -- (sweep);
\draw[arr] (sweep) -- (runner);
\draw[arr] (cli) |- ($(runner.north)+(0,3mm)$) -- (runner.north);
\draw[arr] (runner) -- (registry);
\draw[arr] (runner) -- (result);
\draw[arr] (runner) -- node[lbl, right] {config, seed} (adapter);
\foreach \l in {dm, scale, swarm, mb, cls}
  \draw[thick, {Latex[length=2mm]}-{Latex[length=2mm]}] (adapter.south) -- (\l.north);
\draw[arr] (data.east) -- ++(6mm,0) |- (adapter.west);
\draw[arr] (metrics.east) -| ($(adapter.south west)+(6mm,0)$);
\end{tikzpicture}
\caption{Architecture of the benchmarking platform. A config names a method and a problem; the runner looks the method up in the registry and calls its adapter, which runs the library's own training or evaluation script as a subprocess in that library's directory and normalises what it writes into a \texttt{BenchmarkResult}. Experiment scripts and the sweep driver sit on top as consumers. Every method reads the same seeded dataset files (including the shared starting point $x_0$) and is scored by the same metric functions.}
\label{fig:platform-architecture}
\end{figure}


\subsection{One contract for every method}

A method is anything that implements the abstract class \texttt{L2OMethod}
(\texttt{core/method.py}), which has exactly two operations:
\begin{itemize}
  \item \texttt{train(config, output\_dir) $\to$ TrainOutput}: train the
    optimizer described by \texttt{config} and return a checkpoint handle;
  \item \texttt{evaluate(config, checkpoint, output\_dir) $\to$ EvalOutput}:
    run that checkpoint on the config's problem and return normalised
    metrics.
\end{itemize}
Two kinds of method satisfy this contract. Learned optimizers
(\texttt{trainable = True}) do real work in both calls. Classical optimizers
(Adam, SGD, Adagrad; \texttt{trainable = False}) inherit a no-op
\texttt{train} and implement \texttt{evaluate} by running the classical
update rule through a \emph{host} library's evaluation harness on the same
problem, so that a baseline number is produced by the same code path and
scored the same way as the learned optimizer it is compared with. Methods are
registered by name (\texttt{core/registry.py}); at the time of writing eleven
are registered, eight learned (\texttt{l2o-dm}, \texttt{l2o-rnnprop},
\texttt{l2o-scale}, \texttt{l2o-entropy}, \texttt{l2o-jacobian},
\texttt{l2o-swarm}, \texttt{lista}, \texttt{alista}) and three classical.

The return types are plain data containers (\texttt{core/result.py}). The
one that makes methods comparable is the \emph{trajectory}: a list of scalar
objective values, one per optimization step, from a representative
evaluation. \texttt{summarize\_trajectory} reduces any trajectory to
\texttt{final\_loss}, \texttt{min\_loss}, \texttt{mean\_loss} and
\texttt{num\_points}, and those four numbers appear in every result and in
the leaderboard regardless of which library produced them. The runner
(\texttt{core/runner.py}) drives one benchmark: it resolves the method,
threads the seed (Section~\ref{sec:platform-repro}), calls \texttt{train}
unless asked to evaluate an existing checkpoint, calls \texttt{evaluate},
and writes a \texttt{BenchmarkResult} with the training metrics, evaluation
metrics, trajectory, artifact paths and provenance to \texttt{result.json}.

\subsection{Thin adapters that run the libraries as subprocesses}

The adapter is the translation layer between this contract and one library's
actual scripts. Each adapter builds the library's command-line flags from the
config, runs the library's own \texttt{train.py} or \texttt{evaluate.py}
\emph{as a subprocess in that library's directory}
(\texttt{core/shell.py}), and reads back whatever files the script writes.
Running subprocesses rather than importing the libraries is forced by the
module-name collisions described in Section~\ref{sec:platform-start}, but it
has a second consequence that matters more for a benchmark: the harness never
reaches into a library. The number a method produces inside the platform is
the number its own script produces. The only changes made to the libraries
themselves are the fairness mechanisms of Section~\ref{sec:platform-repro}
(a seed flag, a deterministic evaluation pass, a shared starting point), each
made so that the library's own script yields a number comparable with the
others'.

The adapters differ mainly in what ``one step'' of the trajectory means. For
the model-free optimizers it is one optimizer step on the optimizee. For
LISTA and ALISTA, whose training is a layer-wise Keras \texttt{fit()} rather
than a meta-training loop, the adapter (\texttt{adapters/model\_based.py})
maps each unrolled layer to one trajectory point, so a 16-layer network yields
a 16-point trajectory. The mapping is what lets both families be scored by
the same code, but it must not be read as cost equivalence: one model-free
step runs a recurrent network per coordinate, one unrolled layer is a matrix
product and a threshold. Chapter~\ref{ch:exp2} keeps the two on separate
axes for that reason.

\subsection{Explicit configuration}

A run is described by one JSON config naming the method, the problem, an
optional seed, and two free-form dictionaries \texttt{train} and
\texttt{eval} whose keys are forwarded to the library's scripts as flags.
Two rules make this safe. First, an adapter never falls back to a hidden
default for a field that changes what is trained or evaluated, or how a
checkpoint is located and scored; such keys are looked up with
\texttt{core.config.require}, which raises an error naming the missing key.
The reason is that \texttt{train} and \texttt{evaluate} are separate calls,
and two calls that each default the same field only agree by coincidence.
Second, before a subprocess is launched, \texttt{core/shell.py} scans the
target script's own flag definitions (\texttt{add\_argument} or
\texttt{DEFINE\_*} calls) and rejects any flag the script does not define,
suggesting the closest real one. A mistyped hyperparameter therefore fails
immediately with a message instead of being silently ignored by an
\texttt{argparse} parser several layers down.
